
\documentclass[11pt,a4paper]{article}

\usepackage[utf8]{inputenc}
\usepackage[T1]{fontenc}
\usepackage[french]{babel}
\usepackage[a4paper,margin=2.2cm]{geometry}
\usepackage{booktabs}
\usepackage{graphicx}
\usepackage{float}
\usepackage{hyperref}

\title{\textbf{Étude des performances de la multiplication matricielle parallèle}}
\author{}
\date{\today}

\begin{document}
	
	\maketitle
	
	\section{Introduction}
	
	La multiplication matricielle est une opération fondamentale en calcul
	scientifique. Pour réduire son temps d'exécution, plusieurs techniques de
	parallélisation peuvent être utilisées. Ce travail compare les performances
	d'une version séquentielle, d'une version OpenMP, d'une version MPI et de
	deux configurations hybrides MPI/OpenMP.
	
	L'objectif est d'évaluer l'accélération obtenue et d'identifier les
	configurations les plus adaptées aux tailles de matrices étudiées.
	
	\section{Environnement et méthodologie}
	
	Les expériences ont été réalisées sur une machine équipée d'un processeur
	Intel Core i5-7300U, comprenant 2 cœurs physiques et 4 processeurs logiques,
	avec environ 3,7 Gio de mémoire disponible.
	
	Quatre tailles de matrices carrées ont été testées :
	\[
	N \in \{500,1000,1500,2000\}.
	\]
	
	Les cinq configurations comparées sont :
	
	\begin{itemize}
		\item \textbf{Séquentielle} : 1 processus MPI et 1 thread OpenMP ;
		\item \textbf{OpenMP} : 1 processus MPI et 4 threads OpenMP ;
		\item \textbf{MPI} : 4 processus MPI et 1 thread OpenMP ;
		\item \textbf{Hybride (1$\times$4)} : 1 processus MPI et 4 threads OpenMP ;
		\item \textbf{Hybride (4$\times$1)} : 4 processus MPI et 1 thread OpenMP.
	\end{itemize}
	
	Cinq essais ont été réalisés pour les tailles 500, 1000 et 2000, et dix
	essais pour la taille 1500. Les sommes de contrôle obtenues sont identiques
	pour toutes les configurations d'une même taille, ce qui confirme la
	cohérence numérique des résultats comparés.
	
	\section{Résultats expérimentaux}
	
	L'accélération est définie par :
	\[
	S=\frac{T_{\mathrm{seq}}}{T_{\mathrm{par}}},
	\]
	où $T_{\mathrm{seq}}$ est le temps moyen de calcul séquentiel et
	$T_{\mathrm{par}}$ celui de la configuration parallèle.
	
	L'efficacité est calculée selon :
	\[
	E=\frac{S}{p}\times100\%,
	\]
	où $p$ représente le produit du nombre de processus MPI par le nombre
	de threads OpenMP. Pour la référence séquentielle, l'efficacité de 100\%
	est une convention de présentation.
	
	\begin{table}[H]
		\centering
		\small
		\begin{tabular}{llrrr}
			\toprule
			$N$ & Configuration & Temps (s) & Accélération & Efficacité (\%) \\
			\midrule
			500 & Séquentielle & 0,1590 & 1,000 & 100,00 \\
			500 & OpenMP & 0,0448 & 3,548 & 88,70 \\
			500 & MPI & 0,0607 & 2,622 & 65,56 \\
			500 & Hybride (1$\times$4) & 0,0456 & 3,485 & 87,14 \\
			500 & Hybride (4$\times$1) & 0,0596 & 2,667 & 66,68 \\
			\midrule
			1000 & Séquentielle & 1,3909 & 1,000 & 100,00 \\
			1000 & OpenMP & 0,4153 & 3,349 & 83,74 \\
			1000 & MPI & 0,5542 & 2,510 & 62,74 \\
			1000 & Hybride (1$\times$4) & 0,4038 & 3,444 & 86,11 \\
			1000 & Hybride (4$\times$1) & 0,5494 & 2,532 & 63,29 \\
			\midrule
			1500 & Séquentielle & 4,7705 & 1,000 & 100,00 \\
			1500 & OpenMP & 1,7408 & 2,740 & 68,51 \\
			1500 & MPI & 1,9493 & 2,447 & 61,18 \\
			1500 & Hybride (1$\times$4) & 1,5700 & 3,039 & 75,96 \\
			1500 & Hybride (4$\times$1) & 1,9311 & 2,470 & 61,76 \\
			\midrule
			2000 & Séquentielle & 11,1509 & 1,000 & 100,00 \\
			2000 & OpenMP & 3,8048 & 2,931 & 73,27 \\
			2000 & MPI & 4,5809 & 2,434 & 60,86 \\
			2000 & Hybride (1$\times$4) & 3,9354 & 2,833 & 70,84 \\
			2000 & Hybride (4$\times$1) & 4,7118 & 2,367 & 59,16 \\
			\bottomrule
		\end{tabular}
		\caption{Temps moyens de calcul, accélération et efficacité des configurations.}
	\end{table}
	
	\section{Analyse et discussion}
	
	Les résultats montrent que toutes les configurations parallèles testées
	réduisent le temps de calcul par rapport à la version séquentielle.
	L'accélération observée varie approximativement de 2,4 à 3,5 selon la
	taille des matrices et la configuration.
	
	OpenMP obtient la meilleure accélération pour les matrices de tailles
	500 et 2000. La configuration hybride avec un processus MPI et quatre
	threads OpenMP est la plus rapide pour les tailles 1000 et 1500.
	
	Les configurations reposant sur quatre processus MPI sont généralement
	moins performantes que celles utilisant quatre threads OpenMP. Les coûts
	de communication et de gestion des processus peuvent contribuer à cet
	écart, sans que ces mesures permettent d'isoler précisément chaque cause.
	
	Il faut également tenir compte du matériel : la machine possède deux
	cœurs physiques pour quatre processeurs logiques. Quatre travailleurs
	configurés ne correspondent donc pas nécessairement à quatre cœurs
	physiques indépendants. L'efficacité calculée est un indicateur théorique
	rapporté au nombre de travailleurs configurés, et non une mesure directe
	de l'utilisation du processeur.
	
	Enfin, le temps de calcul doit être distingué du temps total. Le temps
	total inclut des étapes supplémentaires du programme ; les deux mesures
	ne doivent donc pas être interprétées comme équivalentes.
	
	\section{Conclusion}
	
	Cette étude montre que la parallélisation améliore nettement les
	performances de la multiplication matricielle sur la machine utilisée.
	OpenMP et la configuration hybride avec un processus MPI et quatre
	threads OpenMP donnent les meilleurs résultats parmi les configurations
	testées, selon la taille des matrices.
	
	Les résultats dépendent toutefois du matériel, de l'implémentation et
	des coûts associés au parallélisme. Une comparaison sur d'autres
	architectures permettrait de mieux généraliser ces observations.
	
\end{document}
